\documentclass[a4paper,11pt]{article}
\usepackage[italian]{babel}
\usepackage{geometry}
\usepackage{listings}
\geometry{margin=2.5cm}

\title{Verifica Scritta – Basi di Dati (ITS)}
\author{Prof. Fedeli Massimo}
\date{}

\begin{document}
	
	\maketitle
	
	Si consideri il seguente schema logico per la gestione degli stage di un corso ITS.
	
	\bigskip
	
	STUDENTE(\underline{idStudente}, nome, cognome, annoIscrizione, corsoITS)
	
	AZIENDA(\underline{idAzienda}, ragioneSociale, città, settore)
	
	TUTORAZIENDA(\underline{idTutorAz}, nome, cognome, idAzienda)
	
	TUTORACCADEMICO(\underline{idTutorAcc}, nome, cognome, dipartimento)
	
	COMPETENZA(\underline{idCompetenza}, nomeCompetenza)
	
	PROGETTOSTAGE(\underline{idProgetto}, titolo, durataMesi, idAzienda)
	
	STAGE(\underline{idStage}, idStudente, idProgetto, idTutorAz, idTutorAcc, dataInizio, dataFine, valutazione)
	
	STUDENTECOMPETENZA(\underline{idStudente}, \underline{idCompetenza}, livello)
	
	PROGETTOCOMPETENZA(\underline{idProgetto}, \underline{idCompetenza})
	
	\bigskip
	
	\section*{Parte 1 – Query SQL}
	
	Scrivere le seguenti query SQL.
	
	\begin{enumerate}
		
		\item Elencare tutti gli studenti.
		\item Elencare gli studenti iscritti dopo il 2022.
		\item Elencare le aziende del settore `Software''.
		\item Elencare i progetti di stage con durata maggiore di 6 mesi.
		\item Mostrare studenti e azienda presso cui svolgono lo stage.
		\item Elencare gli studenti che hanno svolto almeno uno stage.
		\item Elencare gli studenti che non hanno svolto stage.
		\item Mostrare il numero di stage per ogni azienda.
		\item Calcolare la valutazione media degli stage per ogni studente.
		\item Elencare le aziende che hanno ospitato più di 3 stage.
		\item Mostrare i tutor aziendali con il numero di stage seguiti.
		\item Elencare le competenze possedute dallo studente con idStudente=1.
		\item Elencare i progetti che richiedono la competenza `Java''.
		\item Elencare gli studenti che possiedono la competenza ``SQL''.
		\item Mostrare gli studenti che hanno svolto stage in aziende di città diverse.
		\item Elencare le aziende che non hanno progetti di stage.
		\item Calcolare la durata media dei progetti per azienda.
		\item Elencare gli studenti con valutazione media stage maggiore di 27.
		\item Mostrare il tutor accademico che segue più stage.
		\item Elencare i progetti con il numero di studenti coinvolti.
		\item Mostrare la competenza più richiesta nei progetti.
		\item Elencare gli studenti che possiedono tutte le competenze richieste da un progetto.
		\item Elencare gli studenti che hanno lavorato con più tutor aziendali diversi.
		\item Mostrare la valutazione massima registrata negli stage.
		\item Elencare gli studenti con la media valutazioni più alta.
		
	\end{enumerate}
	
	\section*{Parte 2 – DDL}
	
	Scrivere il codice SQL per creare la tabella:
	
	AULA(idAula, edificio, capienza)
	
	con i seguenti vincoli:
	
	\begin{itemize}
		\item chiave primaria
		\item attributi obbligatori
		\item capienza maggiore di zero
	\end{itemize}
	
	\section*{Parte 3 – Vista}
	
	Creare una vista chiamata:
	
	STUDENTIVALUTAZIONE
	
	contenente:
	
	idStudente, nome, cognome, mediaValutazioneStage.
	
\end{document}
